\clearpage
\VerdierPhysicalPageAnchor{0203}%
\begingroup\small\setlength{\parskip}{0pt}\setlength{\emergencystretch}{1.5em}
\setlength{\abovedisplayskip}{4pt plus 2pt minus 2pt}%
\setlength{\belowdisplayskip}{4pt plus 2pt minus 2pt}%

\VerdierSemanticAnchor{ch3d-2-4-1-cont}
\noindent\textit{est une équivalence de catégories. Le foncteur naturel :}
\[
K^{-}(M^{\wedge})/\operatorname{Ac}^{-}(M^{\wedge})
\longrightarrow D^{-}(\mathcal A)
\]
\noindent\textit{est pleinement fidèle. Les objets de $D^{-}(\mathcal A)$
dont les objets de cohomologie sont des éléments de $M^{\wedge}$ sont
isomorphes à des objets de
$K^{-}(M^{\wedge})/\operatorname{Ac}^{-}(M^{\wedge})$, où
$\operatorname{Ac}^{-}(M)$ et $\operatorname{Ac}^{-}(M^{\wedge})$
désignent les catégories des complexes acycliques dont les composants sont
des éléments de $M$ et $M^{\wedge}$ respectivement.}

\VerdierSemanticAnchor{ch3d-2-4-1-proof}
Cette proposition résulte immédiatement de (2.3.1), (2.3.2) et de
(chap. II, 2.3.5).

\VerdierSemanticAnchor{ch3d-2-4-2}
\medskip\noindent\textbf{Corollaire 2.4.2.}\quad
\textit{Soit $U$ un univers. On suppose que la catégorie $\mathcal A$ est
une $U$-catégorie, i.e.\ que pour tout couple $Y$ et $Z$ d'objets de
$\mathcal A$, $\operatorname{Hom}_{\mathcal A}(Y,Z)$ est un élément de
$U$, et que les sommes directes indexées par des éléments de $U$ sont
représentables dans $\mathcal A$. Soient $(X_i)_{i\in I\in U}$ une famille
de générateurs de $\mathcal A$ et $M$ l'ensemble des objets de $\mathcal A$
de la forme $\coprod_{j\in J\in U}X_j$, où les $X_j$ sont des éléments de
l'ensemble $\bigcup_{i\in I}\{X_i\}$. Alors le foncteur naturel :}
\[
K^{-}(M)/\operatorname{Ac}^{-}(M)\longrightarrow D^{-}(\mathcal A)
\]
\noindent\textit{est une équivalence de catégories.}

\VerdierSemanticAnchor{ch3d-2-4-2-proof}
Résulte immédiatement de (2.4.1).

\VerdierSemanticAnchor{ch3d-2-4-3}
\medskip\noindent\textbf{Proposition 2.4.3.}\quad
\textit{Soient $\mathcal A$ une catégorie abélienne et $M$ un ensemble
d'objets de $\mathcal A$ possédant les propriétés $(E_4)$ et $(E_5)$ tel
que tout objet de $\mathcal A$ soit quotient d'un élément de $M$ et tel
qu'il existe un entier $n\in\mathbb N$ majorant la $M$-dimension de tout
objet $X$ de $\mathcal A$.}

\noindent\textit{a) La catégorie $K^{*}(M)$ est
$\operatorname{Ac}^{*}(\mathcal A)$-localisante à gauche
($*=-,+,b,{}$ ``vide'').}

\noindent\textit{b) Le foncteur naturel :}
\[
K^{*}(M)/\operatorname{Ac}^{*}(M)\longrightarrow D^{*}(\mathcal A)
\]
\noindent\textit{est une équivalence de catégories.}

\VerdierSemanticAnchor{ch3d-2-4-3-proof}
Paraphrase la proposition (2.3.3) ($*=+,b,{}$ ``vide''), ou la proposition
(2.3.1) ($*=-$).

\VerdierSemanticAnchor{ch3d-2-4-4}
\medskip\noindent\textbf{Proposition 2.4.4.}\quad
\textit{On désigne par
$Q^{*}\colon K^{*}(\mathcal A)\longrightarrow D^{*}(\mathcal A)$ le
foncteur canonique.}

\noindent\textit{a) Soit $X$ un objet de $\mathcal A$. L'objet $X$ est
$Q$-libre à gauche (chap. II, 2.3.2) si et seulement si $X$ est un objet
projectif de la catégorie $\mathcal A$.}
\endgroup
